\section{Conclusão}
\label{sec:conclusao}

O trabalho consolidou uma linha metodológica coerente para VRPTW: baseline construtivo por Solomon I1 clássico, seguido de comparação entre VND, GRASP (fixo e reativo) e Busca Tabu sob as mesmas regras de factibilidade e o mesmo objetivo de custo.

Do ponto de vista de engenharia experimental, três resultados são centrais:
\begin{enumerate}
  \item padronização das estruturas de dados e validações, reduzindo risco de inconsistência entre algoritmos;
  \item seleção objetiva de vizinhanças do VND (4 de 8) com base em impacto empírico;
  \item pipeline reprodutível para calibração (\emph{irace}) e geração de tabelas consolidadas por instância.
\end{enumerate}

Com a tabela por instância (gap, custo e $\Delta veh$ por algoritmo), a comparação deixa de ser apenas agregada e passa a mostrar comportamento local por caso, o que é importante para interpretação acadêmica dos resultados em Solomon.

Como próximos passos, recomenda-se:
\begin{enumerate}
  \item concluir o ciclo completo de calibração no orçamento final de \emph{irace};
  \item repetir a campanha consolidada com múltiplas seeds para algoritmos estocásticos;
  \item adicionar testes estatísticos não-paramétricos para comparação de desempenho entre métodos \cite{demsar2006statistical}.
\end{enumerate}
